% Gerado por testes/rodar_testes.py. Não editar à mão.
% Requer \usepackage{longtable}
\subsection*{Autômato de pilha (AP)}
{\footnotesize\setlength{\tabcolsep}{3pt}\begin{longtable}{l p{4cm} r l l p{3.8cm}}
\textbf{ID} & \textbf{Entrada} & \textbf{Tam.} & \textbf{Esperado} & \textbf{Obtido} & \textbf{Regra verificada} \\ \hline \endhead
AP-A1 & \texttt{o,\allowbreak{}a,\allowbreak{}x,\allowbreak{}y\&osrO,\allowbreak{}A,\allowbreak{}X,\allowbreak{}Y\&O.} & 21 & ACEITA & ACEITA & 5 sujeitos e 5 predicativos, pontuação correta (cenário 1) \\
AP-A2 & \texttt{o,\allowbreak{}a,\allowbreak{}x,\allowbreak{}y\&osrO,\allowbreak{}A,\allowbreak{}O,\allowbreak{}Y\&O.} & 21 & ACEITA & ACEITA & Estrutura correta; o AP não verifica a classe (cenário 3) \\
AP-A3 & \texttt{a,\allowbreak{}a,\allowbreak{}a,\allowbreak{}a\&asrA,\allowbreak{}A,\allowbreak{}A,\allowbreak{}A\&A.} & 21 & ACEITA & ACEITA & Todos os itens da mesma classe \\
AP-A4 & \texttt{y,\allowbreak{}x,\allowbreak{}a,\allowbreak{}o\&ysrY,\allowbreak{}X,\allowbreak{}A,\allowbreak{}O\&Y.} & 21 & ACEITA & ACEITA & Classes em outra ordem, contagem igual \\
AP-A5 & \texttt{o,\allowbreak{}o,\allowbreak{}o,\allowbreak{}o,\allowbreak{}o\&osrO,\allowbreak{}O,\allowbreak{}O,\allowbreak{}O,\allowbreak{}O\&O.} & 25 & ACEITA & ACEITA & n = 6: a pilha cresce até 6 S \\
AP-A6 & \texttt{x,\allowbreak{}y,\allowbreak{}x,\allowbreak{}y,\allowbreak{}x\&ysrX,\allowbreak{}Y,\allowbreak{}X,\allowbreak{}Y,\allowbreak{}X\&Y.} & 25 & ACEITA & ACEITA & n = 6 só com plurais \\
AP-A7 & \texttt{o,\allowbreak{}a,\allowbreak{}x,\allowbreak{}y,\allowbreak{}o,\allowbreak{}a\&xsrO,\allowbreak{}A,\allowbreak{}X,\allowbreak{}Y,\allowbreak{}O,\allowbreak{}A\&X.} & 29 & ACEITA & ACEITA & n = 7 \\
AP-A8 & \texttt{a,\allowbreak{}o,\allowbreak{}y,\allowbreak{}x\&osrO,\allowbreak{}O,\allowbreak{}O,\allowbreak{}O\&O.} & 21 & ACEITA & ACEITA & Nenhuma concordância, mas a contagem bate: o AP aceita \\
AP-A9 & \texttt{o\&asrO\&A.} & 9 & ACEITA & ACEITA & Caso-limite: menor cadeia aceita (n = 2) \\
AP-R1 & \texttt{o,\allowbreak{}a,\allowbreak{}x,\allowbreak{}y\&osrO,\allowbreak{}A,\allowbreak{}X\&Y.} & 19 & REJEITA & REJEITA & Falta um predicativo: '.' com S no topo (cenário 2) \\
AP-R2 & \texttt{o,\allowbreak{}a,\allowbreak{}x\&ysrO,\allowbreak{}A,\allowbreak{}X,\allowbreak{}Y,\allowbreak{}O\&A.} & 21 & REJEITA & REJEITA & Predicativos a mais: ',' com $Z_0$ no topo em q7 \\
AP-R3 & \texttt{o,\allowbreak{}a,\allowbreak{}x,\allowbreak{}y,\allowbreak{}osrO,\allowbreak{}A,\allowbreak{}X,\allowbreak{}Y,\allowbreak{}O.} & 21 & REJEITA & REJEITA & Falta o '\&' antes do último item \\
AP-R4 & \texttt{o,\allowbreak{}a,\allowbreak{}x,\allowbreak{}y\&orsO,\allowbreak{}A,\allowbreak{}X,\allowbreak{}Y\&O.} & 21 & REJEITA & REJEITA & 'r' antes de 's' \\
AP-R5 & \texttt{o,\allowbreak{}a,\allowbreak{},\allowbreak{}x,\allowbreak{}y\&osrO,\allowbreak{}A,\allowbreak{}X,\allowbreak{}Y\&O.} & 22 & REJEITA & REJEITA & Duas vírgulas seguidas \\
AP-R6 & \texttt{o,\allowbreak{}a,\allowbreak{}x,\allowbreak{}y\&osrO,\allowbreak{}A,\allowbreak{}X,\allowbreak{}Y\&O} & 20 & REJEITA & REJEITA & Sem o ponto final: termina em q10, não final \\
AP-R7 & \texttt{O,\allowbreak{}A,\allowbreak{}X,\allowbreak{}Y\&OsrO,\allowbreak{}A,\allowbreak{}X,\allowbreak{}Y\&O.} & 21 & REJEITA & REJEITA & Predicativos no lugar dos sujeitos \\
AP-R8 & \texttt{o,\allowbreak{}a,\allowbreak{}x,\allowbreak{}y\&osrO,\allowbreak{}A,\allowbreak{}X,\allowbreak{}Y\&O..} & 22 & REJEITA & REJEITA & Símbolo depois do ponto final \\
AP-R9 & \texttt{(vazia)} & 0 & REJEITA & REJEITA & Caso-limite: entrada vazia (termina em q0) \\
\end{longtable}}
\subsection*{Autômato linearmente limitado (ALL)}
{\footnotesize\setlength{\tabcolsep}{3pt}\begin{longtable}{l p{4cm} r l l p{3.8cm}}
\textbf{ID} & \textbf{Entrada} & \textbf{Tam.} & \textbf{Esperado} & \textbf{Obtido} & \textbf{Regra verificada} \\ \hline \endhead
ALL-A1 & \texttt{oaxyoaxyoOAXYOAXYO} & 18 & ACEITA & ACEITA & 9 pares, todos concordam \\
ALL-A2 & \texttt{ooooooooooOOOOOOOOOO} & 20 & ACEITA & ACEITA & 10 pares masculino singular \\
ALL-A3 & \texttt{aaaaaaaaaAAAAAAAAA} & 18 & ACEITA & ACEITA & 9 pares feminino singular \\
ALL-A4 & \texttt{xyxyxyxyxXYXYXYXYX} & 18 & ACEITA & ACEITA & Só plurais, alternando gênero \\
ALL-A5 & \texttt{yxayoxyaoYXAYOXYAO} & 18 & ACEITA & ACEITA & Classes misturadas, concordância cruzada \\
ALL-A6 & \texttt{oaxyoaxyoaOAXYOAXYOA} & 20 & ACEITA & ACEITA & 10 pares \\
ALL-A7 & \texttt{xxxyyyaaaXXXYYYAAA} & 18 & ACEITA & ACEITA & Blocos de mesma classe \\
ALL-A8 & \texttt{oyoyoyoyoOYOYOYOYO} & 18 & ACEITA & ACEITA & Alterna masculino singular e feminino plural \\
ALL-A9 & \texttt{oO} & 2 & ACEITA & ACEITA & Caso-limite: menor cadeia aceita (1 par) \\
ALL-R1 & \texttt{oaxyoaxyoOAXYOAXYA} & 18 & REJEITA & REJEITA & Último par discorda em gênero (o $\times$ A) \\
ALL-R2 & \texttt{oaxyoaxyoOAOYOAXYO} & 18 & REJEITA & REJEITA & 3º par discorda em número (x $\times$ O) \\
ALL-R3 & \texttt{oaxyoaxyoaOAXYOAXYO} & 19 & REJEITA & REJEITA & Sobra um sujeito (10 $\times$ 9) \\
ALL-R4 & \texttt{oaxyoaxyoOAXYOAXYOO} & 19 & REJEITA & REJEITA & Sobra um predicativo (9 $\times$ 10) \\
ALL-R5 & \texttt{oaxyoaxyoOAXYOAXYOa} & 19 & REJEITA & REJEITA & Sujeito depois dos predicativos \\
ALL-R6 & \texttt{OAXYOAXYOoaxyoaxyo} & 18 & REJEITA & REJEITA & Começa por predicativo \\
ALL-R7 & \texttt{oaxyoaxyoaxyoaxyoa} & 18 & REJEITA & REJEITA & Só sujeitos, nenhum predicativo \\
ALL-R8 & \texttt{oaxyoaxyoOYXAOYXAO} & 18 & REJEITA & REJEITA & Ordem aninhada (espelhada) em vez de cruzada: o 2º par falha \\
ALL-R9 & \texttt{(vazia)} & 0 & REJEITA & REJEITA & Caso-limite: entrada vazia \\
\end{longtable}}
\subsection*{Fluxo completo (AP → ALL)}
{\footnotesize\setlength{\tabcolsep}{3pt}\begin{longtable}{l p{4cm} r l l p{3.8cm}}
\textbf{ID} & \textbf{Entrada} & \textbf{Tam.} & \textbf{Esperado} & \textbf{Obtido} & \textbf{Regra verificada} \\ \hline \endhead
INT-1 & \texttt{o,\allowbreak{}a,\allowbreak{}x,\allowbreak{}y\&osrO,\allowbreak{}A,\allowbreak{}X,\allowbreak{}Y\&O.} & 21 & ACEITA/ACEITA & ACEITA/ACEITA & Estrutura e concordância corretas \\
INT-2 & \texttt{o,\allowbreak{}a,\allowbreak{}x,\allowbreak{}y\&osrO,\allowbreak{}A,\allowbreak{}X\&Y.} & 19 & REJEITA/— & REJEITA/— & 5 sujeitos e 4 predicativos: o ALL não executa \\
INT-3 & \texttt{o,\allowbreak{}a,\allowbreak{}x,\allowbreak{}y\&osrO,\allowbreak{}A,\allowbreak{}O,\allowbreak{}Y\&O.} & 21 & ACEITA/REJEITA & ACEITA/REJEITA & Estrutura correta, 3º par discorda em número (os meninos \ldots{} médico) \\
\end{longtable}}
\subsection*{Símbolo fora do alfabeto}
{\footnotesize\setlength{\tabcolsep}{3pt}\begin{longtable}{l p{4cm} r l l p{3.8cm}}
\textbf{ID} & \textbf{Entrada} & \textbf{Tam.} & \textbf{Esperado} & \textbf{Obtido} & \textbf{Regra verificada} \\ \hline \endhead
ERR-1 & \texttt{o,\allowbreak{}a,\allowbreak{}x,\allowbreak{}b\&osrO,\allowbreak{}A,\allowbreak{}X,\allowbreak{}Y\&O.} & 21 & ERRO & ERRO & Símbolo 'b' fora do alfabeto do AP \\
ERR-2 & \texttt{oaxyoaxyo,\allowbreak{}OAXYOAXYO} & 19 & ERRO & ERRO & Vírgula fora do alfabeto do ALL \\
\end{longtable}}
\paragraph{Mensagens de erro exibidas}
\begin{itemize}
\item \textbf{ERR-1}: \texttt{Erro: o símbolo 'b' (posição 7) não pertence ao alfabeto de entrada \{o, a, x, y, O, A, X, Y, ,, \&, s, r, .\}.}
\item \textbf{ERR-2}: \texttt{Erro: o símbolo ',' (posição 10) não pertence ao alfabeto de entrada \{o, a, x, y, O, A, X, Y\}.}
\end{itemize}
